\section{Introdução}
\label{sec:introducao}
As bombas de fluxo axial constituem uma categoria especializada de turbomáquinas hidráulicas nas quais o movimento do fluido ocorre paralelamente ao eixo de rotação da bomba. Seu princípio de funcionamento baseia-se em um rotor dotado de pás com formato semelhante a hélices de propulsão naval, geralmente contando com 2 a 8 pás que impulsionam o fluido ao longo do eixo. Ao contrário das bombas centrífugas radiais, as bombas axiais mantêm o fluxo na direção axial tanto na entrada quanto na saída.

A característica distintiva das bombas axiais reside em sua velocidade específica elevada, conceito fundamental para classificação de turbomáquinas. A velocidade específica é um parâmetro adimensional calculado pela expressão:
\[
    n_s = \frac{n \sqrt{Q}}{H^{0.75}}
\]
onde $n$ é a rotação (rpm), $Q$ é a vazão (m$^3$/s) e $H$ a altura manométrica (m). Para bombas radiais, $n_s$ tipicamente entre 10--70; diagonais entre 70--120; axiais, acima de 120.


\section{Objetivos}
\label{sec:objetivos}

Este capítulo delineia as metas centrais e as etapas metodológicas fundamentais deste trabalho, estabelecendo as diretrizes para o dimensionamento técnico e a análise operacional da turbomáquina dentro do cenário de estudo proposto.

\subsection{Objetivo Geral}
\label{subsec:obj_geral}

O objetivo principal deste projeto é realizar o dimensionamento termofluidodinâmico, a seleção de materiais e a modelagem cinemática de uma bomba de fluxo axial (modelo Lewis\textsuperscript{\textregistered} VL750) projetada para operar em um sistema de circulação forçada de um evaporador de múltiplo efeito, garantindo a eficiência energética e a confiabilidade operacional frente às severas condições de salinidade, vazão e temperatura do processo.

\subsection{Objetivos Específicos}
\label{subsec:obj_esp}

Para o cumprimento do objetivo geral, as seguintes metas específicas foram estabelecidas e executadas ao longo do estudo:

\begin{itemize}
    \item Dimensionar as grandezas hidráulicas do sistema, incluindo diâmetros de tubulação, fator de atrito e perdas de carga localizadas e distribuídas para o transporte de $8.000~m^3/h$ de solução aquosa salina;
    \item Avaliar a termodinâmica da sucção, calculando o \textit{Net Positive Suction Head} disponível ($NPSH_d$) e propondo soluções de layout de planta (elevação de cota estática) para mitigar o risco de cavitação induzido pela alta pressão de vapor do fluido a $80^\circ C$;
\end{itemize}
